% Figure: adding four phasors of unequal amplitude and phase tip-to-tail in the
% complex plane; the resultant is the closing vector from the first tail to the
% last tip. This is the geometric content of summing several sinusoids at one
% frequency.
\begin{figure}[htbp]
\centering
\begin{tikzpicture}[scale=1.0, line cap=round, >=latex, font=\footnotesize]
  \draw[->, engblack, thin] (-0.6,0) -- (5.6,0) node[right] {$\Re$};
  \draw[->, engblack, thin] (0,-0.6) -- (0,3.4) node[above] {$\Im$};
  % four phasors added tip to tail.
  % p1 = (2.4, 0.0); p2 = (1.0, 1.2); p3 = (-0.4, 1.0); p4 = (1.0, -0.3)
  \coordinate (O) at (0,0);
  \coordinate (P1) at (2.4,0.0);
  \coordinate (P2) at ($(P1)+(1.0,1.2)$);
  \coordinate (P3) at ($(P2)+(-0.4,1.0)$);
  \coordinate (P4) at ($(P3)+(1.0,-0.3)$);
  \draw[engblue, thick, ->] (O) -- (P1) node[midway, below] {$\tilde X_{1}$};
  \draw[engblue, thick, ->] (P1) -- (P2) node[midway, right] {$\tilde X_{2}$};
  \draw[engblue, thick, ->] (P2) -- (P3) node[midway, above] {$\tilde X_{3}$};
  \draw[engblue, thick, ->] (P3) -- (P4) node[midway, above right] {$\tilde X_{4}$};
  % resultant
  \draw[engred, very thick, ->] (O) -- (P4);
  \node[engred, anchor=north west] at ($(P4)+(0.02,-0.05)$) {$\tilde X$};
  \filldraw[engred] (P4) circle (0.05);
  % argument arc of resultant
  \draw[engamber] (0.8,0) arc (0:37:0.8);
  \node[engamber] at (18:1.05) {$\arg\tilde X$};
\end{tikzpicture}
\caption[Adding several phasors tip-to-tail]{Four phasors of unequal amplitude
and phase, summed by the tip-to-tail construction. Each phasor $\tilde X_{n}$
is drawn from the tip of its predecessor; the resultant $\tilde X = \sum_{n}
\tilde X_{n}$ is the single vector from the first tail at the origin to the
last tip. Its magnitude is the amplitude of the combined sinusoid and its
argument is the combined phase. The construction is the same vector addition
that combines two phasors, extended to any number of contributions at a common
frequency; the order of the terms does not change the closing vector, which is
the geometric statement that phasor addition is commutative. When the phasors
are equal in amplitude and evenly spaced in phase the polygon closes on itself
and the resultant is zero, the roots-of-unity cancellation that zeroes a
balanced polyphase current.}
\label{fig:vol02ch03:phasor-polygon}
\end{figure}
